\subsection{引言}\label{sec:7.1}
本章分为两部分.第一部分讨论构造估计量的方法，第二部分讨论如何评价这些（以及其他）估计量.一般而言，这两项活动相互交织.评价估计量的方法往往会启发我们提出新的估计量.不过，暂时我们还是把构造估计量与评价估计量区分开来.

点估计（point estimation）的基本理由很简单.如果从pdf或pmf为$f(x\mid\theta)$的总体抽样，知道$\theta$就意味着知道整个总体.因此，自然会希望找到一个方法，构造点$\theta$的良好估计量，即良好的点估计量.参数$\theta$也常常具有有意义的实际解释（例如总体均值），因而获得$\theta$的良好点估计本身就是我们直接关心的事情.还可能感兴趣的是$\theta$的某个函数，比如$\tau(\theta)$.本章介绍的方法也可用于构造$\tau(\theta)$的估计量.

下面对点估计量的定义看起来也许过于含糊.不过，在这个阶段，我们希望谨慎一些，不把任何候选者排除在考虑范围之外.
\begin{definition}\label{def:7.1.1}\index{SI0763@点估计量}
点估计量（point estimator）是样本的任意函数$W(X_1,\ldots,X_n)$.也就是说，任一统计量都是点估计量.
\end{definition}
注意，这个定义完全没有提到估计量与它要估计的参数之间应有什么对应关系.虽然可以认为定义中应包括这样的说明，但这样会限制可用估计量的集合.定义中也没有提及统计量$W(X_1,\ldots,X_n)$的值域.原则上统计量的值域应与参数的值域一致，但我们将看到情况并非总是如此.

有一点必须区分清楚，即估计值（estimate）与估计量（estimator）的差别.\index{SI0762@点估计值}\semanticindexanchor{762:283}估计量是样本的函数，而估计值是在实际抽取一个样本之后获得的估计量的实现值（即一个数）.在符号上，抽取样本时，估计量是随机变量$X_1,\ldots,X_n$的函数，而估计值是实现值$x_1,\ldots,x_n$的函数.

很多时候，某个参数的点估计量会有显而易见或自然的候选者.例如，样本均值就是总体均值点估计量的自然候选者.不过，一旦离开这种简单情形，直觉不仅可能帮不上忙，甚至可能把我们引入歧途.因此，掌握一些至少能给出合理候选者的技巧很有用.但请注意，这些技巧本身不附带任何保证.它们所产生的点估计量仍须经过评价，才能确立其价值.
